\subsection{Central Simple Algebras}

DEFINITION 1. --- We say that an algebra A over the field K is central if the mapping \(\lambda\mapsto\lambda1\) is a bijection from K to the center of A.

A central algebra is not reduced to 0. For every integer \(n \geqslant 1\) , the matrix K-algebra \(\mathbf{M}_{n}(\mathrm{K})\) is central (VIII, p.~83, Corollary 2) and simple (VIII, p.~120, Theorem 1). More generally, let D be a central K-algebra of finite degree; then \(\mathbf{M}_{n}(\mathrm{D})\) is also central. Let A be a simple ring; its center Z is a field (VIII, p.~121, Corollary 1), and A is therefore a central simple algebra over Z. If the field K is algebraically closed, then a simple algebra of finite degree over K is central (VIII, p.~122, Corollary 3). The opposite algebra of a central simple algebra is central simple.

Remarks. --- 1) Let A and B be K-algebras. If \(A \otimes_{K} B\) is a central simple algebra, then so are A and B (III, §4, No.~4, p.~468, Corollary of Proposition 6 and VIII, p.~221, Proposition 6). Conversely, if A and B are central simple algebras and if one of them has finite degree over K, then \(A \otimes_{K} B\) is a central simple algebra (III, §4, No.~4, p.~468, Corollary of Proposition 6 and VIII, p.~222, Corollary 2).

\begin{enumerate}
\def\labelenumi{\arabic{enumi})}
\setcounter{enumi}{1}
\item
  Let A be a K-algebra and L an extension of the field K. If the L-algebra \(A_{(L)}\) is central simple, then the K-algebra A is central simple. Conversely, if one of the degrees {[}A : K{]} and {[}L : K{]} is finite and if A is a central simple K-algebra, then the L-algebra \(A_{(L)}\) is central simple. This follows from Corollary 2 of VIII, p.~222.
\item
  Let A and B be Morita-equivalent K-algebras. The algebra A is central simple if and only if B is (VIII, p.~105, Proposition 6; p.~111, Corollary; and p.~102, Corollary 1).
\item
  In particular, if A is a central simple K-algebra and \(n \geqslant 1\) , then \(\mathbf{M}_{n}(\mathrm{A})\) is a central simple K-algebra (VIII, p.~102, Example 1).
\end{enumerate}

THEOREM 1. --- Let A be a K-algebra of finite degree. The following properties are equivalent:

\begin{enumerate}
\def\labelenumi{(\roman{enumi})}
\item
  The algebra A is central simple.
\item
  The algebra \(\mathbf{A}\) is central and without radical.
\item
  The canonical homomorphism from the K-algebra \(A \otimes_{K} A^{\circ}\) to the K-algebra \(\operatorname{End}_{K}(A)\) that sends \(a \otimes a'\) to the K-linear mapping \(x \mapsto axa'\) from A to A is bijective.
\item
  There exist an extension L of the field K and an integer \(n \geqslant 1\) such that the L-algebras \(\mathrm{A}_{(\mathrm{L})}\) and \(\mathbf{M}_{n}(\mathrm{L})\) are isomorphic.
\item
  For every separable closure \(\mathrm{K}'\) of \(\mathbf{K}\) , there exists an integer \(n \geqslant 1\) such that the \(\mathrm{K}'\) -algebras \(\mathrm{A}_{(\mathrm{K}')}\) and \(\mathbf{M}_n(\mathrm{K}')\) are isomorphic.
\item
  There exist a Galois extension L of the field K of finite degree and an integer \(n \geqslant 1\) such that the L-algebras \(\mathrm{A}_{(\mathrm{L})}\) and \(\mathbf{M}_n(\mathrm{L})\) are isomorphic.
\item
  There exist a K-algebra of finite degree D that is a field with center K and an integer \(n \geqslant 1\) such that the algebra A is isomorphic to the algebra \(\mathbf{M}_{n}(\mathrm{D})\) .
\end{enumerate}

A ring is simple if and only if it is semisimple and its center is a field (VIII, p.~143, Corollary of Proposition 10). Since A is an algebra of finite degree over the field K, it is a left Artinian ring; it is therefore semisimple if and only if it is without radical (VIII, p.~154, Proposition 4). The equivalence of (i) and (ii) follows.

Set \(E = A \otimes_{K} A^{\circ}\) and \(F = \text{End}_{K}(A)\) ; denote by \(\varphi\) the canonical homomorphism from E to F defined by the relation \(\varphi(a \otimes a')(x) = axa'\) for \(x, a, a'\) in A. If the algebra A is central simple, then so are \(A^{\circ}\) and, therefore, E (Remark 1), so \(\varphi\) is injective. Now, we have \([E : K] = [A : K]^{2} = [F : K]\) , so \(\varphi\) is bijective. Conversely, suppose that \(\varphi\) is bijective; since the algebra F is central simple (because it is isomorphic to a matrix algebra \(\mathbf{M}_{m}(K)\) ), so is E, and therefore so is A (Remark 1). We have thus proved the equivalence of (i) and (iii).

By Remark 4, assertion (vii) implies assertion (i). The converse implication follows from Corollary 3 of VIII, p.~122 and Corollary 2 of VIII, p.~83.

It is clear that (vi) implies (iv), and (iv) implies (i) by Remark 2.

It remains to prove the implications (i)⇒(v)⇒(vi). Suppose that A is central simple, and let \(K'\) be a separable closure of K (V, §7, No.~8, p.~45). Then \(\mathrm{A}_{(\mathrm{K}')}\) is a central simple algebra of finite degree over \(K'\) (VIII, p.~251, Remark 2). By the corollary of VIII, p.~231, there consequently exist an integer \(n \geqslant 1\) and an isomorphism of \(K'\) -algebras from \(\mathrm{A}_{(\mathrm{K}')}\) to \(\mathbf{M}_{n}(\mathrm{K}')\) ; observe that the \(K'\) -algebras \(\mathbf{M}_{n}(\mathrm{K}')\) and \(\mathbf{M}_{n}(\mathrm{K})_{(\mathrm{K}')}\) are isomorphic. By Lemma 4 of VIII, p.~234, there exists a subextension L of \(K'\) , finitely generated over K, such that the L-algebras \(\mathrm{A}_{(\mathrm{L})}\) and \(\mathbf{M}_{n}(\mathrm{K})_{(\mathrm{L})}\) are isomorphic. Then L is separable and of finite degree over K, so contained in a subextension \(L'\) of \(K'\) that is Galois and of finite degree over K (V, §10, No.~1, p.~57, Proposition 2). The \(L'\) -algebras \(\mathrm{A}_{(\mathrm{L}')}\) and \(\mathbf{M}_{n}(\mathrm{L}')\) are then isomorphic.

COROLLARY 1. --- Let A be a central simple algebra of finite degree over a separably closed field K. There exists an integer \(n \geqslant 1\) such that A is isomorphic to the matrix algebra \(\mathbf{M}_{n}(\mathrm{K})\) .

Indeed, every Galois extension of \(\mathrm{K}\) is equal to \(\mathrm{K}\) ; it suffices to apply the equivalence of properties (i) and (v) of Theorem 1.

COROLLARY 2. --- Let A be a central simple algebra of finite degree over K (for example, a field of finite degree over K with center K). There exists an integer \(n \geqslant 1\) such that \([A : K] = n^{2}\) .

Let \(\mathrm{L}\) be an extension of \(\mathrm{K}\) and \(n\) a strictly positive integer such that the L-algebras \(\mathrm{A}_{(\mathrm{L})}\) and \(\mathbf{M}_n(\mathrm{L})\) are isomorphic. We have

\[
[ \mathrm{A}: \mathrm{K} ] = [ \mathrm{A} _ {(\mathrm{L})}: \mathrm{L} ] = [ \mathbf {M} _ {n} (\mathrm{L}): \mathrm{L} ] = n ^ {2}.
\]

In the notation of Corollary 2, the integer n is called the reduced degree of A.

Remark 5. --- Let A be a central simple algebra of finite degree over K whose reduced degree is a prime number \(\ell\) . Then either A is a field, or A is isomorphic to \(\mathbf{M}_{\ell}(\mathrm{K})\) . Indeed, A is isomorphic to an algebra of the form \(\mathbf{M}_{n}(\mathrm{D})\) , where D is a field with center K, and we have

\[
\ell^ {2} = [ \mathrm{A}: \mathrm{K} ] = n ^ {2} [ \mathrm{D}: \mathrm{K} ];
\]

if A is not a field, then \(n \neq 1\) , so \(n = \ell\) and D = K.

